\section{Objetivos}

\begin{frame}{Motivação}
  As Equações Diferenciais Parciais (EDPs) são uma das principais linguagens com que
  modelamos a natureza. Raramente possuem solução fechada, e por isso dependemos de
  \textbf{métodos numéricos} --- tipicamente discretizando o domínio em malhas e avançando
  no tempo.

  \vspace{0.4cm}
  \begin{information}[Uma ponte natural]
    Com a ascensão dos métodos de \emph{Machine Learning} (ML) e de todo um ecossistema de
    \emph{software} construído em torno deles, surge uma extensão natural: aproximar
    soluções de EDPs \alert{aproveitando esse ferramental de Machine Learning} (ML).
  \end{information}

  \vspace{0.3cm}
  Neste trabalho, o objeto de estudo é o \textbf{sistema de Boussinesq}, um modelo clássico
  da hidrodinâmica para ondas longas, não lineares e dispersivas em águas rasas.
\end{frame}

\begin{frame}{Objetivos do Trabalho}
  \begin{block}{Objetivo geral}
    Estudar a \alert{eficácia} de três métodos de Machine Learning --- \textbf{PINN}, \textbf{FNO} e
    \textbf{PINO} --- para aproximar a solução do sistema de Boussinesq, sempre medindo o
    erro contra um \textbf{solver pseudoespectral} de alta acurácia tomado como referência.
  \end{block}

  \vspace{0.2cm}
  Três perguntas guiam o estudo:
  \begin{enumerate}
    \item \textbf{Acurácia vs.\ regime físico:} quão bem cada método reproduz a referência à
          medida que a não linearidade e a dispersão crescem? Onde cada um começa a falhar?
    \item \textbf{Dados, física ou híbrido:} como a presença (ou ausência) de dados
          supervisionados afeta a acurácia e o treinamento? Adicionar um resíduo físico
          ajuda quando há poucos dados?
    \item \textbf{Viés espectral:} como o \emph{spectral bias} se manifesta em cada
          arquitetura, e em que medida dados ou física alteram \emph{quais} frequências do
          erro são resolvidas primeiro?
  \end{enumerate}
\end{frame}
